\documentclass[11pt,a4paper]{article}
\usepackage{fontspec,amsmath,amssymb,amsthm,microtype}
\setmainfont{lmroman10-regular.otf}[BoldFont=lmroman10-bold.otf,ItalicFont=lmroman10-italic.otf,BoldItalicFont=lmroman10-bolditalic.otf]
\usepackage[margin=25mm]{geometry}
\usepackage[unicode,hidelinks]{hyperref}
\hypersetup{pdfauthor={LangevinLean Project}}
\setlength{\emergencystretch}{3em}
\setlength{\parskip}{0.4em}
\setlength{\parindent}{0pt}
\allowdisplaybreaks[2]
\raggedbottom
\title{\texorpdfstring{Compact C2 approximation of the actual Gibbs score and generator}{Compact C2 approximation of the actual Gibbs score and generator}}
\author{LangevinLean Project}
\date{}
\begin{document}\maketitle
Draft mathematical exposition; these notes may contain mistakes.


Smooth radial cutoffs turn an unbounded coordinate of the potential gradient into compact C2 tests. Both the cut-off scores and their true Langevin generators converge in Gibbs L2. The proof derives all score moments from their sub-Gaussian MGFs, controls the product-generator cross terms using exactly the Laplacian-gradient assumption, and applies dominated convergence. This justifies extending compact-test Dynkin identities to the actual score under the paper C3 potential assumption.



Write the genuine scalar generator and vector score generator as

\[g=\nabla f,\qquad \mathcal L\phi=\Delta\phi-\langle g,\nabla\phi\rangle,\qquad G=\nabla\Delta f-Dg\,g.\]

There are compact C2 coordinate tests, with \(\Phi_n=(\phi_{n,j})_{j=1}^p\) and \(\mathcal L\Phi_n=(\mathcal L\phi_{n,j})_{j=1}^p\), for which every displayed vector belongs to Gibbs L2 and

\[\phi_{n,j}\in C_c^2(\mathbb R^p),\qquad \int\!\|\Phi_n-g\|^2\,d\pi\longrightarrow0,\qquad \int\!\|\mathcal L\Phi_n-G\|^2\,d\pi\longrightarrow0.\]
\section{\texorpdfstring{Build an expanding smooth radial cutoff family}{Build an expanding smooth radial cutoff family}}


The nth bump equals one on the ball of radius n+1 and has support inside radius 2(n+1). It is smooth, compactly supported, and takes values in the interval from zero to one. It is exactly one fixed unit-scale bump composed with multiplication by 1/\allowbreak{}(n+1). At every fixed spatial point, all sufficiently large cutoffs are identically one in a neighborhood of that point.


\section{\texorpdfstring{Bound the first and second cutoff derivatives uniformly}{Bound the first and second cutoff derivatives uniformly}}


The fixed compact smooth bump has bounded derivatives through order two. Rescaling by 1/\allowbreak{}(n+1) multiplies its kth derivative by that scalar to the kth power, which is at most one. Hence a common constant bounds the first two derivative norms for every cutoff. Gradient duality converts the first derivative bound into a gradient bound. Expanding the Laplacian in the canonical orthonormal basis bounds its absolute value by p times the same constant.


\section{\texorpdfstring{Derive every finite Gibbs score moment}{Derive every finite Gibbs score moment}}


Each true coordinate of the Gibbs score has the existing Gaussian MGF bound obtained by translation of the density. Sub-Gaussianity supplies all finite scalar moments. Finite-dimensional coordinate membership then proves Lq membership of the entire actual gradient for every finite exponent q. In particular its fourth moment is finite, so one plus its squared norm belongs to Gibbs L2.



Earlier results used here: \cite[Section~3.26.2]{LangevinLean2026ULASC}; \cite[Section~3.26.3]{LangevinLean2026ULASC}; \cite[Section~3.27.3]{LangevinLean2026ULASC}.

\section{\texorpdfstring{Preserve compact C2 regularity and local generator equality}{Preserve compact C2 regularity and local generator equality}}


Multiply a C2 test by the smooth radial cutoff. The product is C2 and has compact support. Because the cutoff is eventually locally one at a fixed point, the products eventually agree with the original test on a neighborhood. Their Frechet derivatives and Laplacians therefore agree there as well. The genuine Langevin generators consequently become pointwise equal. This uses locality of the actual derivatives rather than an assumed approximation identity.



Earlier results used here: \cite[Section~3.46.14]{LangevinLean2026ULASC}; \cite[Section~3.46.2]{LangevinLean2026ULASC}.

\section{\texorpdfstring{Prove dominated L2 convergence for the eventual equalities}{Prove dominated L2 convergence for the eventual equalities}}


If every approximation and its limit are bounded in absolute value by a common L2 function B, their squared differences are bounded by 4 B squared. Eventual pointwise equality makes those differences tend to zero. The ordinary dominated convergence theorem then proves convergence of their actual squared integrals to zero. The module also proves continuity of the actual scalar Langevin generator for C2 tests and the bound M on the gradient norm of every score coordinate.



Earlier results used here: \cite[Section~3.46.14]{LangevinLean2026ULASC}; \cite[Section~3.46.2]{LangevinLean2026ULASC}.

\section{\texorpdfstring{Control the cutoff generator without a bounded third derivative}{Control the cutoff generator without a bounded third derivative}}


The generator product rule gives the cutoff times the original score generator, the score coordinate times the cutoff generator, and twice the inner product of their gradients. The true score-coordinate generator is the corresponding coordinate of gradient(Laplacian f) minus Hessian f applied to gradient f. The first term is bounded by the printed Laplacian-gradient condition and the second by M times the score norm. Uniform cutoff gradient and Laplacian bounds control the remaining terms by a common constant times one plus the squared score norm. This is a genuine Gibbs L2 majorant by the already derived fourth score moment.



Earlier results used here: \cite[Section~3.46.14]{LangevinLean2026ULASC}; \cite[Section~3.46.2]{LangevinLean2026ULASC}.

\section{\texorpdfstring{Produce compact approximations for the real score Dynkin formula}{Produce compact approximations for the real score Dynkin formula}}


For each coordinate j, choose the compact C2 tests equal to the cutoff times the jth actual gradient coordinate. The score cutoffs are dominated by the score norm. Their true generators are dominated by the common L2 majorant just constructed. Both sequences eventually agree pointwise with their targets. Applying the squared-integral dominated convergence principle twice yields convergence in Gibbs L2 of the tests to the coordinate score and of their generators to the corresponding coordinate of the actual gradient generator. Every normalization and moment assumption needed here follows from finite Gibbs partition, Lipschitz gradient, and the source C3 Laplacian-gradient condition.



Earlier results used here: \cite[Section~3.26.2]{LangevinLean2026ULASC}; \cite[Section~3.26.3]{LangevinLean2026ULASC}; \cite[Section~3.26.9]{LangevinLean2026ULASC}; \cite[Section~3.46.2]{LangevinLean2026ULASC}; \cite[Section~3.11.1]{LangevinLean2026IPLA}.

\section{\texorpdfstring{Assemble one vector approximation with explicit L2 membership}{Assemble one vector approximation with explicit L2 membership}}


Choose the proved scalar approximation for each of the finitely many coordinates and place them in the Euclidean coordinate vector. Each compact test and its compact continuous generator is bounded, hence belongs to Gibbs L2. The target generator belongs to Gibbs L2 because its norm is at most M times the true score norm plus the fixed Laplacian-gradient bound. Expanding the squared Euclidean error as a finite sum of squared coordinate errors identifies its integral with the sum of the scalar squared integrals. Finite-sum limit continuity proves convergence of both vector errors to zero. The packaged result includes L2 membership of every vector approximant and both actual targets, so later Cauchy-Schwarz arguments do not assume an unproved moment.



Earlier results used here: \cite[Section~3.26.2]{LangevinLean2026ULASC}; \cite[Section~3.26.3]{LangevinLean2026ULASC}; \cite[Section~3.26.9]{LangevinLean2026ULASC}; \cite[Section~3.46.2]{LangevinLean2026ULASC}.

\begin{thebibliography}{ULA-LSI}
\bibitem[ULA-SC]{LangevinLean2026ULASC}
LangevinLean Project. \emph{Unadjusted Langevin under Strong Convexity}. LL-01, 2026. Draft report, content revision 6f6f50decdb6. \url{https://odakyildiz.github.io/LangevinLean/docs/pdf/UnadjustedLangevin/FullProof.pdf}.
\bibitem[IPLA-SC]{LangevinLean2026IPLA}
LangevinLean Project. \emph{Interacting Particle Langevin under Strong Convexity}. LL-03, 2026. Draft report, content revision c8b5d67ceb43. \url{https://odakyildiz.github.io/LangevinLean/docs/pdf/IPLA/FullProof.pdf}.
\end{thebibliography}

\end{document}
